\documentclass[quick,con]{debate}
\motion{对宠物的爱，爱的更是它 / 爱的更是我}
\stance{反方}{爱的更是我}
\begin{document}
\debatetitle
\begin{thesisbox}[一句话立论]我们按自己的喜好挑了它的长相，按自己的方式表达疼爱，在它的需要和我们的需要真正撞上的时候让路的也是它；所以这份爱最深的那一层，服务的还是我们自己。\end{thesisbox}
\section{定义与判准（标准表述）}
\begin{fields}
\field{爱}{对某个对象有深挚的感情，以及由这份感情带来的持续关切与行动。（这一版我方不争，与正方相同）
}
\field{更}{这是一道比较题，两边都在，比的是分量。我方争的是重心，不是有无。
}
\field{我}{饲主自己的需要，陪伴、被需要、情感的出口、可控感、自我形象。不是“自私”。
}
\field{判准（中立，两条通道）}{一、这份爱的内容是被谁塑造的；二、两边的需要真正撞上时，这份爱跟着谁走。
\begin{bullets}
\item 我方要证明：内容主要由我的缺口与偏好塑造，冲突时被调整的更多是它。
\item 对方要证明：反过来。
\end{bullets}
}
\end{fields}
\section{我方论点}
\begin{longtable}{@{}L{\dimexpr0.041\linewidth-1.50\tabcolsep}L{\dimexpr0.129\linewidth-1.50\tabcolsep}L{\dimexpr0.415\linewidth-1.50\tabcolsep}L{\dimexpr0.415\linewidth-1.50\tabcolsep}@{}}
\toprule \thead{标签} & \thead{一句话} & \thead{核心数据 / 学理 / 案例} & \thead{出处}\\\midrule\endhead
从我的缺口开始 & 从为什么养、养什么、怎么理解它，这份感情的内容是被我的需要和想象塑造的 & 学理：拟人化三因素理论，社交动机越缺人越容易拟人化；拟人化选择把狗选成了像婴儿的脸。数据：285 只犬、68 个品种的研究中，患短头气道综合征的狗的主人 58\% 说狗没有呼吸问题；1359 名英国饲主中自我延伸正向预测积极情绪、负向预测孤独；现实中把它当孩子的 41.7\%，理想中只有 25.4\% & Epley 等 2007, Psychological Review 114(4):864-886；Serpell 2003；Packer 等 2012, Animal Welfare 21(S1):81-93；Ellis 等 2024, Animals 14(3):441；青年志《2025 青年养宠观报告》\\
让路的是它 & 两边需要真正撞上的时候，被调整的是它 & 数据：59\% 的狗与 61\% 的猫超重或肥胖，而 84\% 的狗主人说体型健康；给宠物穿衣热量散不出去，不到一小时可致中暑甚至死亡。案例：中秋节那天送到北京房山救助站的八哥犬“月饼”，因为主人要去外地 & APOP 2022 流行率调查与 2023 观念调查（527 人）；Mota-Rojas 等 2021, Animals 11(11):3263；中国日报网 2025-01-23\\
\bottomrule\end{longtable}
\section{对方论点 → 一句话反驳}
\begin{longtable}{@{}L{\dimexpr0.212\linewidth-1.33\tabcolsep}L{\dimexpr0.381\linewidth-1.33\tabcolsep}L{\dimexpr0.407\linewidth-1.33\tabcolsep}@{}}
\toprule \thead{对方会说} & \thead{我方一句话回} & \thead{对方追问时}\\\midrule\endhead
医疗占 27.6\% & 花钱证明它重要，不证明这份重要是为了谁 & 同一份白皮书里食品占 53.7\%，喂食最容易喂坏它\\
依恋研究：它认得我 & 那是它爱我，这道题问的是我爱谁 & 那些研究的因变量全长在动物身上；它越真实，我方越成立\\
不可替代 & 巴德瓦尔：客观可替换，现象不可替换 & 现象在哪里？在你的感受里\\
武汉封城救援 & 少数人；那五千户占 1.26 亿只的多少 & 托人喂猫成本很低，真正的冲突在弃养名单上\\
照料系统为对方福祉 & 系统是我的，它只是那把钥匙 & Archer：宠物破解了我们育幼反应的密码\\
\bottomrule\end{longtable}
\section{必问与必答}
\begin{fields}
\field{必问 （对方不答就点名）}{
\begin{enumerate}[leftmargin=1.8em,itemsep=2pt,topsep=2pt]
\item 夏天给狗穿裙子拍照，受益人是谁？（一辩稿抛出的那个问题）
\item 你方那些依恋研究，因变量长在谁身上？
\item 同样是花钱，哪一笔算爱它，哪一笔不算？
\item 明知道法斗会喘，为什么它还是最好卖的犬种之一？
\item 为它改变的上限在哪里，会不会辞职？
\end{enumerate}
}
\field{必答 （15 秒以内正面回答）}{
\begin{bullets}
\item “有没有任何一种爱是更指向对方的？” → \textbf{有}。母亲为孩子放弃职业规划。请你方给条件，我们逐条比。
\item “临终的皮下补液怎么解释？” → 止痛的是它，交代是给自己的；而且请问在那之前拖了多久。
\item “你方是不是在说人都自私？” → 不是。我方承认你爱它，我方比的是分量。
\item “肥胖数据是美国的吧？” → 是，这是我方的局限，中国没有可比的体况调查。但服装、写真、美容的比例就在国内的问卷里。
\item “喂胖了只是不懂。” → 不懂是听不见吗？它天天在你面前喘，58\% 的主人还说它没事。
\end{bullets}
}
\end{fields}
\section{自由辩战场概要}
\begin{longtable}{@{}L{\dimexpr0.127\linewidth-1.50\tabcolsep}L{\dimexpr0.289\linewidth-1.50\tabcolsep}L{\dimexpr0.235\linewidth-1.50\tabcolsep}L{\dimexpr0.348\linewidth-1.50\tabcolsep}@{}}
\toprule \thead{战场} & \thead{开场问题} & \thead{目标} & \thead{收束语}\\\midrule\endhead
那件小裙子是给谁穿的 & 夏天给狗穿裙子拍照，受益人是谁？ & 在一个具体画面上让评委看见受益人是人 & 一件它穿不了的衣服，一张它看不懂的照片，受益人写在脸上\\
依恋研究测的是谁 & 那个陌生情境实验，因变量是狗的行为还是人的行为？ & 把正方的核心数据转成我方的证据 & 正方最硬的三组数据，测的全是动物身上的指标\\
边界是谁画的 & 工作调动带不走它，你方怎么办？ & 用生活变动收束，不新开战场 & 那只八哥犬的主人也爱过它，他只是要去外地\\
\bottomrule\end{longtable}
\section{如果……就……}
\begin{contingency}
\ifthen{对方定义不同（把爱定义成“为对方之故”），或拿出没预料到的数据}{当场指出定义杀，问“把狗喂到走不动路算不算爱它”；对数据先问机构、年份、样本，再问它测的因变量在谁身上。}
\ifthen{论点一被打穿（对方用“买之前买之后”把品种论切掉），或对方指控我方犬儒、回避必问}{收缩到论点二，用肥胖与认知落差扛全场，那组数据发生在进门之后；被指控犬儒就撇清“我方承认你爱它”；回避就点名、重复、不换题。}
\end{contingency}
\section{全场三条铁律（背下来）}
\begin{enumerate}[leftmargin=1.8em,itemsep=2pt,topsep=2pt]
\item 每个环节先说一次“我方不否认你爱它”。
\item 被问“有没有更指向对方的爱”必须答“有”，然后要条件。
\item 价值层收在“承认之后才能真的看它一眼”，绝不收在“别骗自己了”。
\end{enumerate}
\end{document}
